\clearpage
\section{Wnioski}

Indeksowa reprezentacja drzewa pozwoliła przechowywać węzły
w jednym wektorze i oddzielić ich rozmieszczenie od logicznych
połączeń. Konsekwencją tego wyboru jest aktualizacja indeksów
po usunięciu pozycji. Nawet gdy węzeł można znaleźć szybko,
udane usunięcie wymaga przejścia po pozostałych elementach
wektora i ma koszt $O(n)$.

Wyszukiwanie wartości i ścieżki oraz porównania podczas
wstawiania zależą od wysokości drzewa $h$. Kolejność danych
ma więc znaczenie. Jeżeli liczby są podawane rosnąco, każdy
następny węzeł staje się prawym dzieckiem poprzedniego.
Wówczas $h$ rośnie liniowo, a budowanie drzewa przez kolejne
wstawienia może wymagać $O(n^2)$ czasu. Dotyczy to również
odtwarzania takiego drzewa z pliku binarnego. Rekurencyjne
wstawianie, inorder, postorder i rysowanie korzystają
ze stosu wywołań o głębokości rzędu $h$.

Każde z trzech przejść odwiedza wszystkie węzły i ma koszt
czasowy $O(n)$, bez uwzględnienia rysunku konsolowego.
Implementacja preorder gromadzi cały wynik w wektorze
i wykorzystuje na ten cel $O(n)$ dodatkowej pamięci.

Próby z podrozdziału~\ref{sec:testy} pokazują, że kolejność
przejścia można dobrać do celu. Inorder daje uporządkowaną
listę kluczy, a preorder wystarcza do zapisania i odtworzenia
kształtu drzewa bez przechowywania indeksów. W obu trybach
odczytu wynik trafia do bieżącego drzewa po sprawdzeniu
danych. Dzięki temu błędny plik nie pozostawia częściowo
wczytanego drzewa.

Naturalnym rozszerzeniem przy wymianie danych między różnymi
platformami byłby format binarny z określoną kolejnością
bajtów i stałą szerokością klucza. Przy pracy z dużymi,
uporządkowanymi zbiorami danych istotnym kierunkiem rozwoju
byłoby ograniczenie wysokości drzewa i głębokości rekurencji.
